%! Counting and Probability with Counting — Unit 3, Lesson 3.5 — Warm-Up
% ONE PAGE, blank and key. Three items at 12pt; do not add a fourth.
% Every number here comes back in the notes:
%   Item 1: 5*4*3*2*1 = 120 (problem 3), 10*9*8 = 720 (problems 4 and 9), 3*2*1 = 6 (row 3, problem 11).
%   Item 2: 7*6*5/(3*2*1) = 35 (problem 7); 10*9*8/(3*2*1) = 120 (problem 10).
%   Item 3: (5/12)(4/11)(3/10) = 60/1320 = 1/22, and 10/220 = 1/22 (problem 15).
% PAGE LOCKSTEP: generate the key with templates/lesson/mkkey.py; every work block is identical.
\documentclass[12pt]{article}
\usepackage{saar-article}
\usepackage{saar-boxes}
\newcommand{\TallMath}[1]{$\displaystyle #1\rule[-1.4em]{0pt}{3.2em}$}
% Word bank strip for every fill-in cluster. Defined per document; the key inherits it and
% prints the same strip, so heights match.
\newcommand{\wordbank}[1]{\par\vspace{2pt}\noindent\fcolorbox{goldacc}{goldbg}{\parbox{\dimexpr\linewidth-2\fboxsep-2\fboxrule\relax}{\small\textbf{Word bank:} #1}}\par\vspace{4pt}}
% Handwriting room inside every \begin{work} block. Moves the blank and the
% key together, so the two can never drift apart.
\setlength{\workrowsep}{30pt}

\begin{document}

\pageheader{Unit 3, Lesson 3.5}{Warm-Up}
% NAMESTRIP: no \namedateperiod here. The student writes their name once, on the
% cover this component is stapled behind.

\vspace{0.15in}

\boxguard
\begin{notesbox}{Spiral Review --- Products and Fractions}
\normalsize
No calculator today. Cancel before you multiply: it keeps the numbers small.

\begin{enumerate}[leftmargin=1.8em, itemsep=18pt, topsep=10pt]

  \item Evaluate each product.
        \wordbank{6 \;\textbullet\; 120 \;\textbullet\; 720}

        \vspace{2pt}
        \renewcommand{\arraystretch}{2.2}
        \begin{tabularx}{\linewidth}{@{} Y Y Y @{}}
        \toprule
        $5 \cdot 4 \cdot 3 \cdot 2 \cdot 1$ & $10 \cdot 9 \cdot 8$ & $3 \cdot 2 \cdot 1$ \\
        \midrule
        \blank{2.4cm} & \blank{2.4cm} & \blank{2.4cm} \\
        \bottomrule
        \end{tabularx}

  \item Cancel common factors first, then multiply.

        \begin{minipage}[t]{0.48\linewidth}
        \begin{work}
          \dfrac{7 \cdot 6 \cdot 5}{3 \cdot 2 \cdot 1} &= 7 \cdot 5 \\
                                                       &= 35
        \end{work}
        \end{minipage}\hfill
        \begin{minipage}[t]{0.48\linewidth}
        \begin{work}
          \dfrac{10 \cdot 9 \cdot 8}{3 \cdot 2 \cdot 1} &= 10 \cdot 3 \cdot 4 \\
                                                        &= 120
        \end{work}
        \end{minipage}

  \item Multiply and simplify. Then simplify the single fraction on the right.

        \begin{minipage}[t]{0.62\linewidth}
        \begin{work}
          \dfrac{5}{12} \cdot \dfrac{4}{11} \cdot \dfrac{3}{10} &= \dfrac{60}{1320} \\
                                                                &= \dfrac{1}{22}
        \end{work}
        \end{minipage}\hfill
        \begin{minipage}[t]{0.34\linewidth}
        \begin{work}
          \dfrac{10}{220} &= \dfrac{1}{22}
        \end{work}
        \end{minipage}

        \vspace{10pt}
        What do you notice about your two answers?
        \par\writelines{1}

\end{enumerate}
\end{notesbox}

\end{document}
